% GENERATED FILE. Edit canonical lessons, then run npm run generate:books.
% canonical-source-hash: fnv1a64:b9f22786cd6d1498
% canonical-lessons: RU-C226-skripka, RU-C226-baraban, RU-C226-tanets, RU-C226-roman, RU-C226-stikhi, RU-R226-first-pass-words-from-chapters-222-223, RU-R226-second-pass-words-from-chapters-224-226

\chapter{Violin, Drum, Dance, Novel, Poems}
\label{ch:ru-violin-drum-dance-novel-poems}
\hlchaptermodality{226}

\begin{chapteropening}
\textbf{By the end of this chapter:} \emph{I can say a violin, a drum, a dance, a novel and poems.}

Complete the last lesson of chapter 226: I can say a violin, a drum, a dance, a novel and poems.
\end{chapteropening}

\section[\texorpdfstring{\ru{скрипка} (skrípka)}{скрипка (skrípka)}]{\ru{скрипка} (skrípka) — a violin}

\label{lesson:RU-C226-skripka}



\begin{quote}
Before the new one: say the Russian for a camel, then the Russian for a deer.
\end{quote}

\subsection*{You'll want to know: \ru{скрипка}}
\textbf{\ru{скрипка}} (\emph{skrípka}) — \textquotedblleft{}a violin\textquotedblright{}.

\begin{grammarlens}[title={What you've built}]
One more word for this chapter.
\end{grammarlens}

\subsection*{Your turn}
\begin{itemize}
\raggedright
  \item \emph{Say it:} \emph{skrípka}
  \item \emph{Say it:} \emph{skrípka}, once more
  \item \emph{Say it:} say \emph{skrípka}
  \item \emph{Recall:} say the Russian for a camel, then the Russian for a deer, then say \emph{skrípka} again
\end{itemize}

\begin{tcolorbox}[breakable,colback=teal!4,colframe=teal!35!black,title={Before you move on}]
What does \ru{скрипка} mean? (\textquotedblleft{}A violin\textquotedblright{}.) Say it once more.
\end{tcolorbox}
\section[\texorpdfstring{\ru{барабан} (barabán)}{барабан (barabán)}]{\ru{барабан} (barabán) — a drum}

\label{lesson:RU-C226-baraban}



\begin{quote}
Before the new one: say the Russian for a deer, then the Russian for a violin.
\end{quote}

\subsection*{You'll want to know: \ru{барабан}}
\textbf{\ru{барабан}} (\emph{barabán}) — \textquotedblleft{}a drum\textquotedblright{}.

\begin{grammarlens}[title={What you've built}]
One more word for this chapter.
\end{grammarlens}

\subsection*{Your turn}
\begin{itemize}
\raggedright
  \item \emph{Say it:} \emph{barabán}
  \item \emph{Say it:} \emph{barabán}, once more
  \item \emph{Say it:} say \emph{barabán}
  \item \emph{Recall:} say the Russian for a deer, then the Russian for a violin, then say \emph{barabán} again
\end{itemize}

\begin{tcolorbox}[breakable,colback=teal!4,colframe=teal!35!black,title={Before you move on}]
What does \ru{барабан} mean? (\textquotedblleft{}A drum\textquotedblright{}.) Say it once more.
\end{tcolorbox}
\section[\texorpdfstring{\ru{танец} (tánets)}{танец (tánets)}]{\ru{танец} (tánets) — a dance}

\label{lesson:RU-C226-tanets}



\begin{quote}
Before the new one: say the Russian for a violin, then the Russian for a drum.
\end{quote}

\subsection*{You'll want to know: \ru{танец}}
\textbf{\ru{танец}} (\emph{tánets}) — \textquotedblleft{}a dance\textquotedblright{}.

\begin{grammarlens}[title={What you've built}]
One more word for this chapter.
\end{grammarlens}

\subsection*{Your turn}
\begin{itemize}
\raggedright
  \item \emph{Say it:} \emph{tánets}
  \item \emph{Say it:} \emph{tánets}, once more
  \item \emph{Say it:} say \emph{tánets}
  \item \emph{Recall:} say the Russian for a violin, then the Russian for a drum, then say \emph{tánets} again
\end{itemize}

\begin{tcolorbox}[breakable,colback=teal!4,colframe=teal!35!black,title={Before you move on}]
What does \ru{танец} mean? (\textquotedblleft{}A dance\textquotedblright{}.) Say it once more.
\end{tcolorbox}
\section[\texorpdfstring{\ru{роман} (román)}{роман (román)}]{\ru{роман} (román) — a novel}

\label{lesson:RU-C226-roman}



\begin{quote}
Before the new one: say the Russian for a drum, then the Russian for a dance.
\end{quote}

\subsection*{You'll want to know: \ru{роман}}
\textbf{\ru{роман}} (\emph{román}) — \textquotedblleft{}a novel\textquotedblright{}.

\begin{grammarlens}[title={What you've built}]
One more word for this chapter.
\end{grammarlens}

\subsection*{Your turn}
\begin{itemize}
\raggedright
  \item \emph{Say it:} \emph{román}
  \item \emph{Say it:} \emph{román}, once more
  \item \emph{Say it:} say \emph{román}
  \item \emph{Recall:} say the Russian for a drum, then the Russian for a dance, then say \emph{román} again
\end{itemize}

\begin{tcolorbox}[breakable,colback=teal!4,colframe=teal!35!black,title={Before you move on}]
What does \ru{роман} mean? (\textquotedblleft{}A novel\textquotedblright{}.) Say it once more.
\end{tcolorbox}
\section[\texorpdfstring{\ru{стихи} (stikhí)}{стихи (stikhí)}]{\ru{стихи} (stikhí) — poems, poetry}

\label{lesson:RU-C226-stikhi}



\begin{quote}
Before the new one: say the Russian for a dance, then the Russian for a novel.
\end{quote}

\subsection*{You'll want to know: \ru{стихи}}
\textbf{\ru{стихи}} (\emph{stikhí}) — \textquotedblleft{}poems, poetry\textquotedblright{}.

\begin{grammarlens}[title={What you've built}]
One more word for this chapter.
\end{grammarlens}

\subsection*{Your turn}
\begin{itemize}
\raggedright
  \item \emph{Say it:} \emph{stikhí}
  \item \emph{Say it:} \emph{stikhí}, once more
  \item \emph{Say it:} say \emph{stikhí}
  \item \emph{Recall:} say the Russian for a dance, then the Russian for a novel, then say \emph{stikhí} again
\end{itemize}

\begin{tcolorbox}[breakable,colback=teal!4,colframe=teal!35!black,title={Before you move on}]
What does \ru{стихи} mean? (\textquotedblleft{}Poems, poetry\textquotedblright{}.) Say it once more.
\end{tcolorbox}
\section[(dialogue)]{Words from chapters 222-223 — a first pass}

\label{lesson:RU-R226-first-pass-words-from-chapters-222-223}



\begin{quote}
Say the words for a novel and poems.
\end{quote}

\subsection*{Your turn}
Say these aloud:

\begin{itemize}
\raggedright
  \item food, a dish, a menu, a bun, kefir
  \item beans, an orange, a lemon, a banana, grapes
\end{itemize}

\begin{tcolorbox}[breakable,colback=teal!4,colframe=teal!35!black,title={Before you move on}]
Food? (\textbf{\ru{еда}}.) A dish? (\textbf{\ru{блюдо}}.) A menu? (\textbf{\ru{меню}}.) A bun? (\textbf{\ru{булка}}.) Kefir? (\textbf{\ru{кефир}}.) Beans? (\textbf{\ru{фасоль}}.) An orange? (\textbf{\ru{апельсин}}.) A lemon? (\textbf{\ru{лимон}}.) A banana? (\textbf{\ru{банан}}.) Grapes? (\textbf{\ru{виноград}}.)
\end{tcolorbox}
\section[(dialogue)]{Words from chapters 224-226 — a second pass}

\label{lesson:RU-R226-second-pass-words-from-chapters-224-226}



\begin{quote}
Say the words for a birch and a spruce.
\end{quote}

\subsection*{Your turn}
Say these aloud:

\begin{itemize}
\raggedright
  \item a birch, a spruce, a pine, a bush, a sheep
  \item a crocodile, a giraffe, a zebra, a camel, a deer
  \item a violin, a drum, a dance, a novel, poems
\end{itemize}

\begin{tcolorbox}[breakable,colback=teal!4,colframe=teal!35!black,title={Before you move on}]
A birch? (\textbf{\ru{берёза}}.) A spruce? (\textbf{\ru{ель}}.) A pine? (\textbf{\ru{сосна}}.) A bush? (\textbf{\ru{куст}}.) A sheep? (\textbf{\ru{овца}}.) A crocodile? (\textbf{\ru{крокодил}}.) A giraffe? (\textbf{\ru{жираф}}.) A zebra? (\textbf{\ru{зебра}}.) A camel? (\textbf{\ru{верблюд}}.) A deer? (\textbf{\ru{олень}}.) A violin? (\textbf{\ru{скрипка}}.) A drum? (\textbf{\ru{барабан}}.) A dance? (\textbf{\ru{танец}}.) A novel? (\textbf{\ru{роман}}.) Poems? (\textbf{\ru{стихи}}.)
\end{tcolorbox}
